\section{Project Diary}

The full GitHub commit log is the canonical record of what was done and when; what follows is a curated narrative grouped into five phases, with representative commits referenced by short SHA so that the reader can drill into the diff online. Author initials are used for brevity: \textbf{LP}~=~Leonardo Palanca, \textbf{SP}~=~Sebastian Pasz, \textbf{LM}~=~Lorenzo Mottinelli, \textbf{JM}~=~Johnny Magon.

\subsection{Phase 1 -- Bootstrap (April 18, 2026)}
The repository was created with a placeholder README (\texttt{c27c04b}, LP) and four initial pushes from each team member, mostly establishing folder layout, a Python virtual environment and dependency list (\texttt{1dacea2}, LP), and a first SQLite \texttt{test\_db.py} (\texttt{9999da4}, LP). The first non-trivial component was a hidden-folder scraper exploration (\texttt{9a9def8}, ``don't look here scrapers'', LP) which fed into a \texttt{MonitorAgent} prototype the same evening (\texttt{39d344f}, JM). A first HTML dashboard mockup and an architecture sketch landed alongside (\texttt{128816b}, SP), followed by improvements to the monitor (\texttt{0a8731e}, JM) and a working \texttt{fetch\_doc\_text} helper (\texttt{8989d1a}, JM). The choice to use Anthropic-style commit messages and to keep the database URL secret (\texttt{80dbfb0}, LP) was made on day one.

\subsection{Phase 2 -- Architecture Consolidation (April 30, 2026)}
After a pause to clarify the brief, the team rewrote the AGENTS contract and made it the entry point for any contributor or AI agent (\texttt{0abc90c}, LP). The README was consolidated and the architecture section moved into it (\texttt{299a76a}, LP). A synthetic Fed-like static site, \texttt{FakeFed}, was deployed behind Nginx as a reproducible scraper fixture (\texttt{b738e84}, LP), and the homepage of \texttt{fedwatcher.ellep.it} was published as a static landing dashboard (\texttt{1779a51}, LP). Dead files from the bootstrap phase were removed (\texttt{d1935b0}, \texttt{73c51ce}, LP).

\subsection{Phase 3 -- Analyst v1 and Macro Ingestion (May 5--13, 2026)}
The first end-to-end \texttt{AnalystAgent} appeared on May 5: it segments each FOMC document into rhetorical sections (assessment, outlook, policy stance, vote) and asks an LLM to fill a structured rubric per section (\texttt{b8370d2}, SP). Documentation followed in lockstep (\texttt{a48a688}, \texttt{8f8fb75}, \texttt{2b38c26}, LP and SP). OpenRouter was added as a vendor-agnostic LLM gateway so the team could swap between Claude and OpenAI without code changes (\texttt{6ff8127}, LP). On May 7 the FRED backfill landed for CPI, core CPI, unemployment and 2y yields (commit grouped under several README updates by LP, see \texttt{4a1014e}). Mid-May the monitor was made source-aware so it can ingest both the live Fed site and \texttt{FakeFed} (\texttt{b209be1}, LP), and macro refresh was extended to multiple series (\texttt{f8bf297}, LP).

\subsection{Phase 4 -- Multi-Analyst, Strategist and Automation (May 19--22, 2026)}
The team intentionally produced several competing analyst implementations to compare segmentation strategies, prompt styles and LLM backends. ``Better analyst'' and ``better analyst 2'' (\texttt{3b93a26}, \texttt{34c1243}, SP) explored alternative section weights, and a \emph{dual-model} analyst was added that runs two LLMs in parallel and reconciles their outputs (\texttt{bcaf99b}, JM; database wiring in \texttt{c1e1b49}, SP). The \texttt{StrategistAgent} got its first real version: divergence between tone-implied and market-implied next-rate was flipped to the economically correct sign (\texttt{f327b1f}, LM) and extended with explicit narrative output (\texttt{72ec4d9}, LM). Cron-driven automation was added in three pieces: a bash wrapper layer (\texttt{8ccbc66}, \texttt{e10e03f}, SP), a 90-day documents update window (\texttt{904f796}, SP), and a final ``last agent'' that closes the loop (\texttt{e62cdfe}, JM). The historical scraper that backfills FOMC pages all the way to 1994 was moved into \texttt{scripts/} (\texttt{278e98c}, LM). A weighted-score bug -- zero-weighted sections were polluting the average -- was fixed shortly afterward (\texttt{b70b358}, JM).

\subsection{Phase 5 -- Productionization and Ops (May 22--27, 2026)}
The static \texttt{snapshot.json} that initially powered the dashboard was replaced by a live read-only FastAPI service (\texttt{bc8cd76}, LP). The dashboard's data explorer received a legend and labels to make table semantics obvious (\texttt{32c0c38}, LP). The FOMC \texttt{FEDFUNDS} policy rate was backfilled into \texttt{macro\_data} to unblock the strategist (\texttt{7a3c04c}, LM). A new \texttt{report/} directory replaced the older \texttt{literature.md} file (\texttt{39f75bd}, LM). The schema was updated to match the new analyst variants (\texttt{f53ecbf}, JM) and the README was rewritten to document all four agents end-to-end (\texttt{e1e0afd}, JM).

The final commit at the time of writing, \texttt{6cbaa7a} (LP, May 27), tightens \texttt{AGENTS.md} so that any contributor -- human or AI -- must treat \texttt{db/schema.sql} as authoritative and must explicitly flag missing tables or columns rather than silently introducing them. This is a direct consequence of an operational incident on the same day, where two SQLite copies of the production database had drifted: the API was reading a complete \texttt{/opt/FEDWatcher/fedwatcher.db} while all cron writers were updating an incomplete \texttt{/home/programming/FEDWatcher/fedwatcher.db}. The two databases were merged, the live API was repointed to the cron-written copy, and the redundant \texttt{/opt} tree was archived and removed. The schema-contract commit closes the door on similar drift in the future.

\subsection{Tools and Workflow}
Across all phases the team relied on a small, stable toolchain: \textbf{Python 3} (FastAPI, \texttt{requests}, \texttt{beautifulsoup4}, \texttt{pandas}, \texttt{sqlite3}), \textbf{SQLite} for storage, \textbf{Nginx} and \texttt{systemd} on the VPS, and \textbf{vanilla HTML/CSS/JS} (Chart.js) for the dashboard, deliberately avoiding a frontend framework. Generative-AI tools are listed in Section~\ref{sec:ai}.
